\section{Methods}

This section outlines the research design, model adoption, process of developing the system, system evaluation, sampling population, data collection, data processing, and ethical considerations are discussed. Ensuring the steps taken for this study are well-defined to answer the objectives of the study. Furthermore, this also ensures that the study follows a clear structure that can be followed and replicated by researchers in the future.

\subsection{Research Design}

The research design for this study follows Applied and Quantitative research design. Applied aspect is for the adoption and implementation of mathematical models for itinerary generation and route optimization. The models that are adopted for this study are the Adaptive Genetic Algorithm with Dynamic Mutation and Crossover Probabilities (AGAM) for personalized itinerary generation, and various metaheuristic algorithms for solving the Traveling Salesman Provincial (TSP). The Quantitative is for the evaluation of the developed system using the Technology Acceptance Model (TAM) and ISO/IEC 25010:2023 standards.

\subsection{Model Adoption}

This section discusses the optimization problems such as route optimization and itinerary generation and the steps in modeling each mathematically. Specifically, the mathematical models used in the study are the traveling salesman problem (TSP) and the orienteering problem (OP).

The following metaheuristic algorithms for route optimization were considered for this study. This is due to the nature of metaheuristic algorithms to not guarantee an optimal solution, but rather a near-optimal solution within a reasonable time frame. These algorithm includes:

\begin{enumerate}
	
\item \textbf{Genetic Algorithm (GA).} Inspired by evolution and natural selection, the Genetic Algorithm is an optimization technique that works by analyzing potential solutions and generating new ones through the processes of selection, crossover, and mutation. This iterative process continues until a predefined termination requirement is met. GA is widely used across various domains because of its ability to navigate complex search spaces and identify optimal or near-optimal solutions \parencite{Hossain2024}.

\item \textbf{Ant Colony Optimization (ACO).} The Ant Colony Optimization (ACO) is an optimization algorithm that mimics the foraging behavior of ants. It utilizes artificial pheromone trails to guide the search for optimal solutions. Ants in the algorithm uses probabilistic pheromones and heuristic information to address problems. The quality of the solutions determines how the pheromone trails are updated, this mechanism encourages exploration and prevents premature convergence \parencite{Hossain2024}.

\item \textbf{Simulated Annealing (SA).} The Simulated Annealing (SA) algorithm is a metaheuristic algorithm based on the annealing process in metallurgy. It begins with an initial solution and incorporates a mechanism that allows for the occasional acceptance of a ``worse'' solution based on current temperature and more thoroughly explore the solution space. The probability of accepting a worse solution is gradually reduced over time, simulating the cooling process \parencite{Hossain2024}.

\item \textbf{Particle Swarm Optimization (PSO).} The Particle Swarm Optimization (PSO) algorithm is inspired by the social behavior of bird flocks and fish schools. In PSO, the population is made up of numerous particles, each representing a potential solution within the search space. Every particle is characterized by a velocity and a fitness value. Conceptually, PSO is very straightforward, as it relies solely on elementary mathematical operations and does not require any derived information about the optimization function \parencite{Wadi2025}.

\item \textbf{Elephant Herding Optimization (EHO).} The Elephant Herding Optimization (EHO) algorithm is a swarm-based metaheuristic search approach for addressing optimization problems that mimics the real herding behavior of elephants, which involves organizing elephant swarms into smaller clans, each led by a matriarch and consisting of multiple female elephants and their calves, while adult male calves are expelled from the clan \parencite{Wadi2025}.

\item \textbf{Gray Wolf Optimizer (GWO).} The Gray Wolf Optimization (GWO) algorithm is a metaheuristic technique inspired by gray wolves, simulating their leadership and cooperative hunting behaviors to locate optimal solutions. The hierarchical structure of a wolf pack (Alpha, Beta, Delta, and Omega wolves) is shown in the algorithm’s design. Wolves’ positions are updated based on the Alpha’s location and the calculated hunting distance. GWO leverages a balance of exploration and exploitation to discover the best possible solutions \parencite{Hossain2024}. 

\item \textbf{Hovering Scouts and Foraging Flocks Pied Kingfisher Optimizer (HSFFPKO).} HSFFPKO enhances the traditional Pied Kingfisher Optimizer by introducing multiscale search mechanisms to improve optimization performance. By adding Hovering Scouts for fine-grained local adjustments and Foraging Flocks for parallel group-based exploitation, the algorithm maintains better population diversity than its predecessor. These additions allow HSFFPKO to mitigate local optima entrapment and solve complex, high-dimensional problems more efficiently \parencite{DelaCruz2025HSFFPKO}.
\end{enumerate}

To determine the best algorithm to implement for Lakad, a comparative analysis was conducted on several metaheuristic algorithms used for solving TSP. The analysis focuses on solution quality (accuracy) and runtime speed of each algorithm. The algorithms were tested on selected instances from the TSPLIB library, a standard benchmark for TSP and related problems. Table \ref{tsplib_instace} presents selected instances for benchmarking. 

\begin{table}[H]
	\centering
	\singlespacing
	\caption{TSPLIB Instances for Benchmarking}
	\label{tsplib_instace} \renewcommand{\arraystretch}{1.2}
	\begin{tabularx}{\linewidth}{XX >{\centering\arraybackslash}X}
		\toprule
		\textbf{TSPLIB instance} & \textbf{Description} & \textbf{Optimal Solution} \\
		\midrule
		burma14 & Small instance with 14 cities & 3323\\ 
		bays29 & Medium instance with 29 cities & 2020\\ 
		eil51 & Large instance with 51 cities & 426\\ 
		berlin52 & Large instance with 52 cities & 7542\\
		\bottomrule
	\end{tabularx}
\end{table}

The experiments were implemented in Python because of its library support for TSPLIB, and for simplicity of the language. To ensure statistical reliability and mitigate the stochastic nature of heuristic algorithms, each algorithm was executed for 10 independent iterations on each TSPLIB instance. This was based on studies of \textcite{Hossain2024,Wadi2025}, where repeated independent runs are used to account for the stochastic nature of metaheuristics algorithms, as a single run may be differed from other runs. The results were averaged to get the baseline performance profile. Two primary metrics were defined to evaluate the algorithms:

\begin{enumerate}
	\item \textbf{Runtime Speed (\( T_{avg} \)).} It measures the average time it takes to find the solution across the 10 iterations and measured in seconds.  The calculation is shown in Equation \ref{eq:runtime}: 
	\begin{equation}
		\label{eq:runtime}
		T_{avg} = \frac{1}{10} \sum_{i=1}^{10} T_i
	\end{equation}
	Where $T _{i}$ is the runtime for the $i$-th iteration.

\item \textbf{Solution Quality (\( Q_{avg} \)).} The accuracy of algorithm is measured with average relative error since the optimal solution for each TSPLIB instance is known. This is computed in Equation \ref{eq:quality}
	\begin{equation}
		\label{eq:quality}
		Q_{avg} = \frac{1}{10} \sum_{i=1}^{10} \frac{C_i - C_{opt}}{C_{opt}} \times 100
	\end{equation}
	Where $C_i$ is the cost of the solution found in the $i$-th iteration, and $C_{opt}$ is the cost of the known optimal solution for that TSPLIB instance.
\end{enumerate}

All algorithms were executed under identical testing conditions. The tests were performed on a computer system that has Ryzen 5 3400G CPU and 16 gigabytes of RAM running Arch Linux operating system. The researchers developed the algorithms through Jupyter Notebook development environment and they used Python built-in time module for runtime measurement.

To select the final algorithm, the trade-off between speed and accuracy was assessed using Simple Additive Weighting (SAW). The researchers chose it because of its simplicity and interpretability. Allowing for a straightforward comparison of the algorithms based on the defined metrics. Since both $T _{avg}$ and $Q_{avg}$ are “cost” criteria, meaning a lower values represent better performance, it is normalized with respect to the maximum (worst) observed value to scale all metrics between 0 and 1. The normalization formula used was shown in Equation \ref{eq:norm}:
\begin{equation}
	\label{eq:norm}
	r _{ij} = \frac{x_{ij}}{\max(x_{ij})}
\end{equation}
Where $r _{ij}$ is the normalized value of the $j$-th criteria for the $i$-th algorithm, $x_{ij}$ is the original value of the $j$-th criteria for the $i$-th algorithm, and $\max(x_{ij})$ is the maximum value observed for the $j$-th criteria across all algorithms. This normalization ensures that both metrics are on a comparable scale, allowing for a fair assessment of each algorithm’s performance in terms of both speed and accuracy. 

The final score $S _{i}$ for each algorithm $i$ was calculated in Equation \ref{eq:saw}: 
\begin{equation}
	\label{eq:saw}
	S _{i} = \sum_{j=1}^{n} w_j \cdot r _{ij}
\end{equation}
Where $S _{i}$ is the final score for the $i$-th algorithm, $w_j$ is the weight assigned to the $j$-th criteria, such that $\sum_{j=1}^{n} w_j = 1$, and $r _{ij}$ is the normalized value of the $j$-th criteria for the $i$-th algorithm.

For this study, since the weights can be subjectively assigned, equal weights were assigned to both criteria, i.e., $w _{1} = w _{2} = 0.5$. This gives equal importance of runtime speed and solution quality, the final score is shown in Equation \ref{eq:finalscore}. And the algorithm with the lowest score $S_{i}$ is considered the best choice for implementation for Lakad.
\begin{equation}
	\label{eq:finalscore}
	S _{i} = 0.5 \cdot r _{i1} + 0.5 \cdot r _{i2}
\end{equation}


Another algorithm that is adopted for this study was the modified variant of the genetic algorithm developed by \textcite{Yochum2020} called Adaptive Genetic Algorithm with Dynamic Mutation and Crossover Probabilities (AGAM) for the personalized itinerary generation. AGAM addresses itinerary planning as a Multi-Objective Optimization Problem, an optimization technique rooting from orienteering problems. As discussed in their work, \textcite{Yochum2020} developed AGAM to generate optimal and personalized travel itineraries by balancing several objectives such as point of interest (POI) rating, and distance of POIs between the next one. Rooting from Genetic Algorithm (GA), AGAM simulates the process of natural evolution and utilizes genetic operators such as selection, crossover, and mutation to iteratively improve the quality of itineraries which, unlike traditional GA that use fixed mutation rates, also employs dynamic adjustments in parameters during the evolutionary process in a way that can maintain population diversity and therefore generate efficient and adaptive solutions. Even though AGAM can also be used for route optimization, the study chose to use it for itinerary generation because of its effectiveness in generating personalized itineraries as shown in the study of \textcite{Yochum2020}.

In this implementation, the population set is defined as $P = \left\{ p _{1},\dots, p _{n} \right\}$, where each individual $p _{i}$ represents a potential itinerary. An itinerary is encoded as a sequence of POIs, $p _{i}= \left\{ q _{1}, q _{2}, \dots, q _{m} \right\}$, where each gene $q _{i}$ corresponds to a specific POI visit. The length of the sequence depends on the number of POIs selected, and it is constrained by a user-defined limit $MAXN$. The total travel distance of the itinerary is constrained by $MAXD$, which ensures that the generated itineraries does not exceed the user’s travel distance budget. 

The quality of each itinerary is evaluated by the fitness function given in Equation \ref{agamFitness}. Since factors such as popularity, rating, and interest are on different scales, a normalization process is applied to each factor so that its values is between 0 and 1. 

\begin{equation}
	\label{agamFitness}
	f(p_i) = w_{1}Ti(p_i) + w_{2}Tn(p_i) + w_{3}Tr(p_i) 
\end{equation}
Where $w _{j}$ for $j = 1, 2, 3$, represents the weight assigned to each factor, and it is defined as follows:
\begin{itemize}
    
    \item $Ti(p_i)$ is the sum of all interest values of each POI in the itinerary $p_i$. Each POI contains an interest value which quantifies how interested the user is in that specific POI category. The formula for calculating $Ti(p_i)$ is shown in Equation \ref{Ti},
    \begin{equation}
		\label{Ti}
		Ti ( p _{i} ) = \sum _{q_i \in p_i} I(q_i)
	\end{equation}
	where $I(q_i)$ is a binary value that indicates whether the user is interested in the category of POI $q_i$ (1 if interested, 0 otherwise).
    
    \item $Tn (p _{i})$ is the total number of POIs in $p _{i}$, which is calculated in Equation \ref{Tn}.
    \begin{equation}
        \label{Tn}
        Tn ( p _{i}) = \frac{N _{op} ( p _{i} )}{NOP}, \quad \text{where } N _{op} ( p _{i} ) \le MAXN 
    \end{equation}
    where $N _{op} ( p _{i} )$ is the number of POIs in itinerary $p _{i}$, and $NOP$ is the total number of available POIs.
    
    \item $Tr( p _{i} )$ is the total rating of $p_i$ given by Equation \ref{Tr},
    \begin{equation}
        \label{Tr}
        Tr ( p _{i} ) = \frac{\sum_{q_i \in p_i} R(q_i)}{MAXR}
    \end{equation} 
    where $R(q_i)$ is the rating of a specific POI $q_i$ in itinerary $p_i$ obtained from Google Maps, and $MAXR$ is the best possible rating of all POIs defined as $MAXR = \text{Number of POIs} \times \text{Maximum Rating}$.
\end{itemize}

AGAM incorporates dynamic adjustments in mutation and crossover probabilities to balance exploration (searching new solutions) and exploitation (refining existing solutions). These are defined as follows:
\begin{itemize}
    \item \textbf{Adaptive Crossover $ ( P _{c} )$.} To preserve high-quality solutions, the crossover probability decreases as the fitness of the parents exceeds the population average $ ( f _{avg} )$. This is given by Equation \ref{crossover},
    \begin{equation}
        \label{crossover}
        PC = \begin{cases}
            pc_1 - \dfrac{(pc_1 - pc_2)(f' - f_{avg})}{f_{max} - f_{avg}}, & f' \ge f_{avg} \\
            pc_1, & f' < f_{avg}
        \end{cases}
    \end{equation}
    Here, $pc_1$ and $pc_2$ are the upper and lower bounds for crossover probability, $f_{max}$ is the largest valued fitness score among the set of populations $P$, $f'$ is the larger fitness score of the two parents being crossed, and $f_{avg}$ refers to the average fitness score of the population $P$. 

    \item \textbf{Adaptive Mutation $ ( P _{m} )$.} Similarly, mutation rates are adjusted to prevent destroying optimal solutions. This is defined in Equation \ref{mutation},
    \begin{equation}
        \label{mutation}
        PM = \begin{cases}
            pm_1 - \dfrac{(pm_1 - pm_2)(f - f_{avg})}{f_{max} - f_{avg}}, & f \ge f_{avg} \\
            pm_1, & f < f_{avg}
        \end{cases}
    \end{equation}
    where $pm_1$ and $pm_2$  are the bounds for mutation probability, and $f$ is the fitness score of the individual being mutated.
\end{itemize}

For the parameters of AGAM, the following values were used: $pc_1 = 0.9$, $pc_2 = 0.6$, $pm_1 = 0.1$, and $pm_2 = 0.01$. These values were chosen based on \textcite{Yochum2020} recommendation, which are commonly used in dynamic crossover and mutation probabilities. 

\subsection{Process of Developing the System}

The development of the Lakad mobile application uses the AGILE methodology, with a specific focus on the Kanban method. This approach integrates all key software development phases. These phase include analysis, planning, design, development, testing, and review, each of the phases contributes to the overall quality and functionality of the final product.

\begin{figure}[H]
	\centering
	\caption{Software Development Life Cycle}
	\label{kanban}
	 \includegraphics[width=\textwidth, trim={0 6cm 0 4cm},clip]{SDLC.png}
\end{figure}

Figure \ref{kanban} shows a graphical representation of a Kanban workflow model.Kanban method is use because of its effective visual process management. It optimize the workflow and allowing the researchers to track the progress of developement. \textcite{Alaidaros2021} highlighted that Kanban operates on the Just In Time principle, meaning tasks are completed and delivered while maximizing efficiency. Kanban’s adaptability allows the team to adjust priorities and reallocate resources as needed. This flexibility is key to maintaining momentum and delivering a high-quality software output \parencite{Risener2022}.


\begin{figure}[H]
	\centering
	\caption{System Architecture}
	\label{SysArch}
	\includegraphics[width=0.8\textwidth]{SysArch.png}
\end{figure}

Figure \ref{SysArch} shows the system architecture of Lakad shows how various application components work together to produce optimize travel routes for users. The mobile frontend functions as the system’s primary interface which enables tourists to search for destinations and create itineraries while they organize their travel arrangements. The application frontend is connected to Mapbox API to handle user interactions through which users request map data and location coordinates. The API processes these requests and returns the corresponding responses, which are then displayed on the user interface as maps and travel routes. Supabase enables the application to store user-created itineraries together with additional data in an external database which users can use to view and control their itineraries. The admin users use their dedicated admin GUI which keeps them separate from tourists to handle user-visible location content while they calculate distance matrices using OpenRouteService matrix API.

\begin{figure}[H]
	\centering
	\caption{Entity Relationship Diagram}
	\label{erd}
	\includegraphics[width=0.8\textwidth]{erd.jpeg}
\end{figure}

The Entity Relationship Diagram (ERD) of the system, as shown in Figure \ref{erd}, consists of eight entities; \texttt{place}, \texttt{opening\_hours}, \texttt{distances}, \texttt{stops}, \texttt{itinerary}, \texttt{profiles}, \texttt{reviews}, and \texttt{review\_reports}. A \texttt{profiles} entity can generate many itineraries of the entity \texttt{itinerary} and provide multiple entries in the \texttt{reviews} entity. An \texttt{itinerary} entity can contain as many stops, and acts as a junction entity for a place before being organized inside an itinerary, which lets the system properly assign a \texttt{visit\_order} and \texttt{visit\_duration} for a location within a specific trip. \texttt{stops} entities reference a place, which serves as the central container for location data, coordinates, and ratings. Additionally, each place has at least one dedicated \texttt{opening\_hours} entry to define its schedule. A place entity can be the subject of multiple reviews. A \texttt{reviews} entity may also have its own \texttt{review\_reports} which are then used by the system to track moderation status.

\begin{figure}[H]
	\centering
	\caption{Data Flow Diagram}
	\label{dfd}
	\includegraphics[width=0.8\textwidth]{DFD.png}
\end{figure}

Two types of users exist for the whole system which is the tourist and admins as illustrated in Figure \ref{dfd}. There are two main flows for the tourist actions: place flow and itinerary flow. The  tourists can view and add places to their itineraries. They can also perform create, read, update, and delete (CRUD) operations on their itineraries. One part of the itinerary flow starts from the user generating an itinerary through the AGAM process is sent to the itinerary management process for the visualization through graphical user interface and is saved to the database. Tourists can also optimize the stops in the itineraries and request navigation instructions for the itinerary. Admins on the other hand, manages what the users can see. They can perform CRUD operations on the places that are seen by the tourists.

\begin{table}[H]
	\centering
	\singlespacing
	\caption{List of Tools for Development}
	\label{listtools}
	\renewcommand{\arraystretch}{1.2}
	\begin{tabularx}{\linewidth}{p{0.3\textwidth}X}
		\toprule 
		\textbf{Name} & \textbf{Description} \\
		\midrule
		Visual Studio Code  & A powerful text editor developed by Microsoft. It supports multiple programming languages and offers built-in debugging, version control integration, and extensions suitable for React Native development. \\
		React Native        & An open-source framework developed by Meta for building cross-platform mobile applications using JavaScript and React. \\
		JavaScript     & An Interpreted, high-level programming language for building user interfaces, interaction and application logic.\\
		Git                 & A software for versioning and managing software projects. \\
		GitHub              & Host for git repositories that allows online collaboration. \\
		\bottomrule
	\end{tabularx}
\end{table}

Table \ref{listtools} presents the list of tools used for the development. Lakad was developed using the React Native main development framework, which uses fast and cost-efficient cross-platform mobile development, allowing a single codebase to function on both iOS and Android devices. The system utilized the JavaScript programming language for implementing the application’s user interface and core functionalities due to its flexibility and compatibility with React Native framework. The main development environment for writing, editing, and debugging the program was Visual Studio Code, which provided lightweight performance and had a large selection of extensions appropriate for React Native development. Git and GitHub was used throughout the development process to maintain organized team work and track software versions.

\subsection{System Evaluation}

The target population for this study consists of 150 end-users and 22 IT experts. The end-users consists of tourists who visiting a tourist spot in selected municipalities of Bulacan, namely Baliwag, Bocaue, Bustos, Calumpit, Guiguinto, Malolos, Marilao, Plaridel, Pulilan, and San Rafael. The 150 end-users was then distributed proportionately across the municipalities based on the tourists arival data from then Provincial Historical and Cultural Tourism Office (PHACTO) of Bulacan. For the IT experts, 22 IT experts where deliberately selected based on their professional qualifications and relevant experience. IT consists of but not limited to software engineers, UI/UX designers, Quality Assurance professionals and faculty members from The College of Information and Communications Technology (CICT) at Bulacan State University (BulSU).

\subsubsection{Data Collection}

The data collection process of this study was divided into two parts, one for the end-users (tourists) and another for the IT professionals. Each was evaluated using different instruments to gather relevant data for assessing the system from both user experience and technical quality perspectives. The end-users were evaluated using TAM, while the IT professionals were evaluated using ISO/IEC 25010:2023. The TAM questionnaire for end-users uses 5-point Likert scale to assess the acceptability of Lakad. For the ISO/IEC 25010:2023 evaluation for IT professionals uses 5-point numerical scale. These scales were adopted from existing instruments and were modified to fit the context of this study especially for the TAM questionnaire. 

The data collection process began after the local government units of the selected municipalities granted their formal authorization to conduct tourist surveys. A pilot testing was conducted after receiving all the required permissions to assess the TAM questionnaire’s reliability and validity. The pilot test involved a group of respondents who shared characteristics with the target study population. The pilot test selected 20 respondents from municipality of Pandi because this group would not participate in the main study to prevent any potential research bias. The research team used Cronbach’s alpha to assess the internal consistency and reliability of the adapted TAM questionnaire based on pilot test results. The computed Cronbach’s alpha coefficient of 0.881 confirmed that all TAM items demonstrated good internal consistency while maintaining reliable performance for user technology acceptance measurement within this research study.

The data collection for end-users was conducted through face-to-face surveys in various tourist spots across the selected municipalities in Bulacan via convenience sampling. The researchers approached tourists along with the informed consent and invited them to participate in the survey, which included a brief introduction to Lakad and its features. After exploring the application they were asked to complete the TAM questionnaire. Additionally, demographic information such as age bracket, travelling frequency, and technological proficiency was collected to analyze potential correlations between these factors and the TAM constructs. 

For the ISO/IEC 25010:2023, a permissions to conduct the survey among IT professionals was requested and obtained from the Information and Communications Technology Office (ICTO) of BulSU and to the dean of CICT. The data collection for IT professionals was conducted through an online survey or in-person sessions, depending on the availability and preference of the respondents. The IT professionals were provided with access to Lakad and were given a brief overview of its features and functionalities. After exploring the application, they were asked to complete the ISO/IEC 25010:2023 evaluation. In addition to the quantitative evaluation, IT professionals were also encouraged to provide qualitative feedback on the system’s strengths and areas for improvement based on their technical expertise.  

\subsubsection{Data Processing and Analysis}

A systematic data processing and analysis was done for this study. The collected data from TAM questionnaires and ISO/IEC 25010:2023 evaluations were analyzed and interpreted according to the results. The TAM questionnaire utilized a five-point Likert-type scale to measure respondents’ perception towards Lakad. The instrument consists of declarative statements that is designed to measure the four TAM constructs. For the ISO/IEC 25010:2023 evaluation instrument used a five-point performance rating scale to evaluate system quality constructs through expert assessment. The encoding for the scales of both instruments are presented in Table \ref{tab:scale_encoding}.

\begin{table}[H]
	\centering
	\singlespacing
	\caption{Encoding of Measurement Scales for TAM and ISO/IEC 25010:2023}
	\label{tab:scale_encoding}
	\renewcommand{\arraystretch}{1.2}
	\begin{tabularx}{\linewidth}{>{\centering\arraybackslash}X >{\centering\arraybackslash}X >{\centering\arraybackslash}X}
		\toprule
		\textbf{Scale/Code} & \textbf{TAM} & \textbf{ISO/IEC 25010:2023} \\ 
		\midrule
		1      & Strongly disagree & Poor       \\
		2      & Disagree          & Fair       \\
		3      & Neutral           & Good       \\
		4      & Agree             & Very good  \\
		5      & Strongly agree    & Excellent  \\ 
		\bottomrule
	\end{tabularx}
\end{table}

The responses from individual questionnaire items are ordinal in nature, descriptive statistics were utilized to summarize the data without assuming equal intervals between response categories. At the item level, frequency distributions and median values were computed to describe response patterns and determine the central tendency of each item. To obtain construct-level measures, composite scores were computed by averaging responses across items corresponding to each construct. For both TAM and ISO/IEC 25010:2023 instruments, the composite score for the $i$-th respondent was calculated using in Equation \ref{Xi},

\begin{equation}
	\label{Xi}
	X_i = \frac{\sum_{j=1}^{n} Y_{ij}}{n}
\end{equation}
Where, $X_{i}$ is the composite score for the $i$-th respondent, $Y_{ij}$ is the rating of the $i$-th respondent on the $j$-evaluation item, and $n$ is the total number of item to evaluated for a given construct. 

The overall mean score for each construct was computed by taking the mean of the composite scores of all respondents, this is given by Equation \ref{Xbar},

\begin{equation}
	\label{Xbar}
	\bar{X} = \frac{\sum_{i=1}^{m} X_i}{m}
\end{equation}
Where, $\bar{X}$ is the overall mean score for a given construct, $X_i$ is the composite score for the $i$-th respondent, and $m$ is the total number of respondents. The mean provided a summary measure of overall user acceptance (for TAM constructs) and overall system quality (for ISO/IEC 25010:2023 constructs). 

The sample standard deviation was also calculated to get the variability of response among the evaluators. A lower standard deviation indicates that the respondents were in strong agreement regarding Lakad’s system acceptance and quality.  The formula for sample standard deviation is given by Equation \ref{S},

\begin{equation}
	\label{S}
	s = \sqrt{\frac{\sum_{i=1}^{m} (X_i - \bar{X})^2}{m-1}}
\end{equation}
Where, $s$ is the sample standard deviation, $X_i$ is the composite score for the $i$-th respondent, $\bar{X}$ is the overall mean score for a given construct, and $m$ is the total number of respondents.

\subsection{Ethical Considerations}
